\documentclass[11pt,a4paper]{letter}
\usepackage[margin=2.5cm]{geometry}
\signature{Maniru Ibrahim\\Mathematics Applications Consortium for Science and Industry (MACSI)\\Department of Mathematics and Statistics, University of Limerick, Ireland\\Maniru.Ibrahim@ul.ie}
\address{}
\date{\today}
\begin{document}
\begin{letter}{The Editor-in-Chief\\Socio-Economic Planning Sciences}
\opening{Dear Editor,}

We submit the manuscript ``A conformally calibrated linear-fractional robust model for food-indexed cash transfers and price monitoring under partial information'' by Maniru Ibrahim, Garba Abor Isa and Ibrahim Abdullahi Shehu for consideration as a research article in \emph{Socio-Economic Planning Sciences}.

Cash transfers are often paid at one national rate even when food prices differ widely across regions. Indexing payments to local prices requires data that planners rarely have: statistical agencies publish national or zonal averages, the most and least expensive regions, and surveys that arrive late. In Nigeria, state food price tables stopped after August 2023, a few months after the removal of the petrol subsidy. The paper asks how a fixed budget should be allocated across regions under such partial information, how much protection that allocation can guarantee, and which regions are worth monitoring.

The paper makes three contributions. First, it shows that the allocation with the best worst-case efficiency has a closed form, pays each region in proportion to its largest possible need share, and can be computed exactly by linear programming, by sorting with a Dinkelbach iteration, or in closed form, depending on the information available. Second, it calibrates the uncertainty set by conformal prediction, so the guarantee holds with a stated probability, and extends the calibration to data that become stale. Third, it uses the guarantee to choose which regions to monitor and shows that reporting delay limits the value of monitoring more than the number of monitored regions.

The application uses 125,598 state--item--month food prices for Nigeria and the food budgets of poor households from the 2018/19 living standards survey. Over 2021 to 2023 a flat transfer delivered 78\% of the protection the same budget could have delivered, while a guarantee-constrained rule delivered 91\% and certified 64\% in the worst case. We believe the paper fits the journal's interest in quantitative methods for public-sector planning and social policy, and that it will interest readers working on social protection, humanitarian operations and official statistics.

The manuscript is original, has not been published, and is not under consideration elsewhere. All authors have approved the submission. The authors declare no competing interests, and the research received no specific funding. Code and derived data are available at https://github.com/Manirmaths/robust-food-indexed-transfers.

\closing{Yours sincerely,}
\end{letter}
\end{document}
